% !TeX root = closed-systems-from-comparison-completeness.tex
% ============================================================
\chapter{Transport Closure and the Uniqueness of the Macroscopic Compatibility Law}
\label{ch:einstein}
% ============================================================
\begin{chapterorientation}
\orientationitem{Input}{The quadratic carrier, causal scaling, Hilbert scalar channel, and second-jet faithful realized branch.}
\orientationitem{Role}{This chapter classifies the admissible second-order macroscopic compatibility law.}
\orientationitem{Output}{The Einstein-boundary law as a compatibility statement for realized geometry, used downstream by connection-first reconstruction and phase curvature.}
\end{chapterorientation}
\section{Introduction}
Under the standing principle of closed-world admissibility
(\cref{sp:closed-world}), this chapter develops the macroscopic
compatibility consequences of the intrinsic transport package
already established in \cref{ch:interface,ch:causality,ch:hilbert}.
Using the stabilized quadratic carrier together with the inherited
observer-side filtration and, under the inherited second-jet
faithfulness condition, the realized degree--\(2\) channel and
first-variation pairing carried there, it derives the unique
macroscopic compatibility law used in \cref{ch:connection,ch:em}.
Curvature and Einstein-tensor notation in the realized branch follows
the standard general-relativistic convention \cite{Wald1984}.
\Cref{ch:interface,ch:causality,ch:hilbert} therefore supply the
exact intrinsic transport, filtration, and realized-branch input
needed here, rather than a bare whole-carrier curvature
identification or an unconditional direct smooth-realization
package.
At the intrinsic level, the first nontrivial stabilized obstruction is
the quadratic transport carrier
\[
	F^2/F^3
\]
on the Stone inverse limit.
At the realized level, two structural facts have already been established.
First, \cref{ch:interface} shows that, under the inherited
second-jet faithfulness condition, nonzero stabilized square classes
are equivalent to nonzero realized curvature in the realized metric
geometry.
Second, the pointwise first-variation geometry constructed in
Proposition~17.2.2, together with the smooth realization theorem for
connection-first geometry in Theorem~17.9.1(ii) and the tangent
identification stated in Section~17.12, identifies the intrinsic
first-variation arenas with the realized tangent geometry.  More
precisely, for each realized point \(p\),
\[
	V_p \cong T_pM,
\]
and, under the inherited second-jet faithfulness condition, the
first-variation pairing on the realized degree--\(2\) channel
attached to the stabilized quadratic carrier
\[
	B_p:V_p\times V_p\to\mathbb R,
\]
is identified with the Lorentz metric pairing
\[
	g_p:T_pM\times T_pM\to\mathbb R.
\]
Thus the realized transport geometry contains not only the curvature
channel of the quadratic carrier, but also, under that inherited
second-jet faithfulness condition, the pointwise quadratic pairing
carried by the realized degree--\(2\) channel attached to the
stabilized quadratic carrier.
In particular, comparison automorphisms preserve the realized transport
geometry at the metric level and not merely at the curvature level.
The purpose of the present chapter is to determine which macroscopic
compatibility laws are admissible for that realized metric geometry.
No action functional is assumed.
No matter tensor is postulated.
No variational principle, entropy extremization principle, or external
geometric axiom is inserted.
The argument is entirely internal to the closed stack and uses only
structures already determined by
\cref{thm:interface-main,thm:lorentz-forced,thm:quadratic-scalar,thm:universal-complex-structure}.
We work strictly after the development in
\cref{ch:interface,ch:causality,ch:hilbert}.
In particular, we take as input the following previously established
facts.
\begin{enumerate}[label=\textup{(S\arabic*)}]
	\item \label{ein:S1}
	      Rectangular completeness and canonical quotient semantics.
	\item \label{ein:S2}
	      Closed-system descent and exhaustion of enrichment by the same
	      two loci, possibly with both and with no third: representative
	      choice and transport.
	\item \label{ein:S3}
	      Transport closure:
	      every admissible microscopic dynamics lies in
	      \[
		      \G:=\Aut(U,\mathcal C).
	      \]
	\item \label{ein:S4}
	      A refinement tower
	      \[
		      \B_1\subset \B_2\subset \cdots,
		      \qquad
		      \B_\infty=\Bigl\langle \bigcup_k \B_k \Bigr\rangle,
	      \]
	      with Stone inverse limit
	      \[
		      S_\infty=\UF(\B_\infty)\cong \varprojlim_k \UF(\B_k).
	      \]
	\item \label{ein:S5}
	      Foundation's quantifier boundary:
	      global admissibility is equivalent to finite non-coverage of the
	      excluded finite comparison cells, equivalently to existence of a
	      coherent admissible ultrafilter tower.
	\item \label{ein:S6}
	      The stabilized intrinsic degree--\(2\) transport carrier
	      \[
		      \mathcal K_\infty\subset F^2/F^3
	      \]
	      on the inverse-limit locus.
	\item \label{ein:S7}
	      Interface realization:
	      under the inherited second-jet faithfulness condition and smooth
	      realization
	      \[
		      \iota:S_\infty\to M,
	      \]
	      nonzero stabilized square classes are equivalent to nonzero
	      realized curvature for a smooth realized metric \(g\).
	\item \label{ein:S8}
	      First-variation realization:
	      by Proposition~17.2.2, Theorem~17.9.1(ii), and Section~17.12,
	      the pointwise first-variation space \(V_p\) realizes as the
	      tangent space \(T_pM\), and, under the inherited second-jet faithfulness condition, the first-variation pairing on the
	      realized degree--\(2\) channel attached to the stabilized
	      quadratic carrier
	      \[
		      B_p:V_p\times V_p\to\mathbb R
	      \]
	      is identified with the Lorentz metric pairing
	      \[
		      g_p:T_pM\times T_pM\to\mathbb R.
	      \]
	\item \label{ein:S9}
	      Local rigidity:
	      every homeomorphism of \(M\) preserving the realized transport geometry
	      is a \(C^1\)-diffeomorphism.
	\item \label{ein:S10}
	      Quadratic determinacy and causal-diamond growth:
	      along the observer-side filtration inherited from
	      \cref{ch:flow} and interleaved with the refinement tower in
	      \cref{ch:stitching},
	      \[
		      \mu(D_T)\asymp T^2,
	      \]
	      and equivalently, after reparameterizing by the corresponding
	      refinement scale \(L:=k(T)\),
	      \[
		      \mu(D_L)\asymp L^4.
	      \]
\end{enumerate}
The question is therefore the following.
\begin{quote}
	Which macroscopic compatibility laws are admissible for the realized
	metric geometry determined by the closed relational stack?
\end{quote}
The answer is that, in the metric-only regime of the closed stack, the
only admissible law is the Einstein law with cosmological term.
\begin{theorem}[Macroscopic compatibility classification]
	\label{thm:einstein-main}
	Under the standing principle \cref{sp:closed-world}, let
	\[
		E:\mathcal G\to \Gamma(S^2T^*M)
	\]
	be an admissible macroscopic operator in the sense of
	\cref{def:admissible-operator}.
	Then there exist constants \(\alpha,\beta\in\mathbb R\) such that
	\[
		E_{\mu\nu}(g)=\alpha\,G_{\mu\nu}(g)+\beta\,g_{\mu\nu}.
	\]
	Consequently, every admissible macroscopic compatibility law
	\[
		E(g)=0
	\]
	is equivalent, when \(\alpha\neq 0\), to
	\[
		G_{\mu\nu}(g)+\Lambda g_{\mu\nu}=0,
		\qquad
		\Lambda:=\frac{\beta}{\alpha}.
	\]
\end{theorem}
The proof proceeds in five steps.
\begin{enumerate}[label=\textup{(\roman*)}]
	\item finite comparison-preserving partial symmetries extend globally;
	\item nontrivial automorphisms can be localized to arbitrary clopen
	      cylinders in the Stone inverse limit;
	\item under realization, those localized symmetries preserve the full
	      realized transport geometry, including, under the inherited
	      second-jet faithfulness condition, the pointwise first-variation
	      pairing carried by the realized degree--\(2\) channel attached
	      to the stabilized quadratic carrier, and hence induce local
	      \(C^1\)-diffeomorphisms of the realized manifold preserving the
	      metric pairings \(g_p\);
	\item finite comparison data at a realized point reconstructs the
	      quadratic form on the tangent space, and the induced tangent action
	      realizes the full orthogonal group \(O(T_pM,g_p)\);
	\item in dimension \(4\), every admissible local second-order symmetric
	      metric operator is a linear combination of \(G\) and \(g\), and the
	      divergence identity then follows from the contracted Bianchi identity.
\end{enumerate}
We then strengthen the conclusion by showing that the cosmological term
is not an external order--\(0\) residue.
It is the scalar projection of the realized degree--\(2\) channel
attached to the stabilized quadratic carrier, read under the
inherited second-jet faithfulness condition through the same
intrinsic obstruction mechanism that supports the Einstein branch.
% ============================================================
\section{Stack boundary and admissible macroscopic operators}
\label{sec:einstein-boundary}
% ============================================================
\begin{definition}[Admissible macroscopic operator]
	\label{def:admissible-operator}
	Let \(\mathcal G\) denote the realized metric class on \(M\).
	A map
	\[
		E:\mathcal G\to \Gamma(S^2T^*M)
	\]
	is called \emph{admissible} if the following properties hold.
	\begin{enumerate}[label=\textup{(A\arabic*)}]
		\item \label{ein:A1}
		      \textbf{Locality of order at most \(2\).}
		      For every \(x\in M\), the value \(E(g)(x)\) depends only on the
		      \(2\)-jet \(j_x^2g\).
		\item \label{ein:A2}
		      \textbf{Metric-only dependence.}
		      The operator depends only on the realized metric \(g\).
		\item \label{ein:A3}
		      \textbf{Symmetry.}
		      For every \(g\in\mathcal G\),
		      \[
			      E_{\mu\nu}(g)=E_{\nu\mu}(g).
		      \]
		\item \label{ein:A4}
		      \textbf{Realized local symmetry invariance.}
		      Whenever \(\Phi\) is a realized local symmetry arising from a
		      cylinder-supported automorphism of \(S_\infty\), one has
		      \[
			      E(\Phi^*g)=\Phi^*E(g).
		      \]
		\item \label{ein:A5}
		      \textbf{Concrete second-order tensoriality.}
		      In local coordinates, each component \(E_{\mu\nu}(g)\) is a finite sum
		      of complete contractions built from
		      \[
			      g_{\alpha\beta},
			      \qquad
			      (g^{-1})^{\alpha\beta},
			      \qquad
			      \partial_\gamma g_{\alpha\beta},
			      \qquad
			      \partial_{\gamma\delta} g_{\alpha\beta},
		      \]
		      with smooth coefficient functions.
	\end{enumerate}
\end{definition}
\begin{remark}
	Condition \textup{\ref{ein:A4}} supplies the microscopic symmetry input
	determined by the stack.
	The classification argument below does not require density in the full
	diffeomorphism group.
	What it requires is the tangent isotropy determined by realized local
	comparison symmetries.
\end{remark}
\begin{remark}[Local algebraicity of admissible operators]
	By \textup{\ref{ein:A1}}, \textup{\ref{ein:A2}}, and
	\textup{\ref{ein:A5}}, the value \(E(g)(p)\) depends only on the
	\(2\)-jet of \(g\) at \(p\), and after passing to normal coordinates at
	\(p\), that dependence is an algebraic construction from the algebraic
	curvature tensor and the metric.
	Thus once the orthogonal isotropy at \(p\) is identified, the
	classification reduces to ordinary orthogonal invariant theory.
\end{remark}
% ============================================================
\section{Finite typing and extension from the quantifier boundary}
\label{sec:einstein-extension}
% ============================================================
\begin{definition}[Finite comparison constraints]
	Let \(F\subset U\) be finite.
	A finite directed comparison pattern over \(F\) is a specification of
	bits
	\[
		(\pi^\rightarrow_{c,f},\pi^\leftarrow_{c,f})\in\{0,1\}\times\{0,1\}
		\qquad
		(c\in\mathcal C,\ f\in F).
	\]
	Its realization set is
	\[
		X(F;\pi^\rightarrow,\pi^\leftarrow)
		:=
		\Bigl\{
		u\in U:
		c(u,f)=\pi^\rightarrow_{c,f},\
		c(f,u)=\pi^\leftarrow_{c,f}
		\text{ for all }(c,f)
		\Bigr\}.
	\]
\end{definition}
\begin{lemma}[Typing lemma]
	\label{lem:einstein-typing}
	For every finite set \(F\subset U\) and every finite pattern
	\((\pi^\rightarrow,\pi^\leftarrow)\), the realization set
	\[
		X(F;\pi^\rightarrow,\pi^\leftarrow)
	\]
	belongs to some finite-stage Boolean algebra \(\B_k\).
\end{lemma}
\begin{proof}
	Each atomic condition \(c(u,f)=1\), \(c(u,f)=0\), \(c(f,u)=1\), or
	\(c(f,u)=0\) defines a comparison half-space.
	The realization set of the finite pattern is therefore a finite Boolean
	combination of finitely many such half-spaces.
	By construction of the refinement tower, every finitely generated
	Boolean subalgebra of \(\B_\infty\) is contained in some finite stage
	\(\B_k\).
	Hence
	\[
		X(F;\pi^\rightarrow,\pi^\leftarrow)\in \B_k
	\]
	for some \(k\).
\end{proof}
\begin{theorem}[Extension from the quantifier boundary]
	\label{thm:einstein-extension}
	Let
	\[
		\sigma:D\to D'
	\]
	be a finite partial bijection such that
	\[
		c(x,y)=c(\sigma(x),\sigma(y))
		\qquad
		(x,y\in D,\ c\in\mathcal C).
	\]
	Then \(\sigma\) extends to an element
	\[
		\phi\in \G=\Aut(U,\mathcal C).
	\]
\end{theorem}
\begin{proof}
	Fix a well-order of \(U\), and perform a back-and-forth construction.
	Assume inductively that a finite comparison-preserving partial bijection
	\[
		\sigma_n:D_n\to D_n'
	\]
	has already been constructed.
	\smallskip
	\noindent
	\emph{Forth step.}
	Choose \(u\in U\setminus D_n\).
	We seek \(v\in U\setminus D_n'\) such that for every \(x\in D_n\) and
	every \(c\in\mathcal C\),
	\[
		c(u,x)=c(v,\sigma_n(x)),
		\qquad
		c(x,u)=c(\sigma_n(x),v).
	\]
	This prescribes a finite comparison pattern over the witness set
	\(D_n'\).
	By \cref{lem:einstein-typing}, the set of realizations of this pattern
	lies in some \(\B_k\).
	Suppose this realization set were empty.
	Then the finitely many excluded comparison cells determined by the
	required pattern would cover \(U\), contradicting the quantifier
	boundary assumption \ref{ein:S5}.
	Hence the realization set is nonempty.
	Because only finitely many points have already been used, we may choose
	\(v\notin D_n'\).
	\smallskip
	\noindent
	\emph{Back step.}
	Apply the same argument to the inverse partial bijection in order to
	extend the codomain side by one point.
	\smallskip
	Proceeding transfinitely along the chosen well-order yields a total
	bijection
	\[
		\phi:U\to U
	\]
	preserving every comparison in \(\mathcal C\).
	Hence
	\[
		\phi\in\Aut(U,\mathcal C)=\G.
	\]
\end{proof}
% ============================================================
\section{Cylinder-local automorphism completeness on the Stone limit}
\label{sec:einstein-cylinder}
% ============================================================
\begin{definition}[Stone cylinders]
	For \(b\in \B_\infty\), define
	\[
		\widehat b:=\{p\in S_\infty:\ b\in p\}.
	\]
	A \emph{clopen cylinder} is a set of this form.
\end{definition}
\begin{theorem}[Cylinder-local completion]
	\label{thm:einstein-cylinder}
	For every nonempty clopen cylinder \(E\subseteq S_\infty\), there
	exists a nontrivial automorphism \(\phi\in\G\) such that the induced
	Stone homeomorphism \(\widehat\phi\) satisfies
	\[
		\widehat\phi|_{S_\infty\setminus E}=\id.
	\]
\end{theorem}
\begin{proof}
	Write \(E=\widehat b\) for some \(b\in \B_\infty\).
	Because \(b\) belongs to the direct-limit Boolean algebra, it is
	decided by finitely many generators from some stage \(\B_k\).
	Let \(W\) be a finite witness family supporting those generators.
	Choose an ultrafilter \(p\in E\).
	Since \(E\) is nonempty and clopen, and since finite-stage refinement
	separates finite comparison protocols, there exists another ultrafilter
	\(q\in E\), \(q\neq p\), which agrees with \(p\) on every generator
	supported outside \(W\), but differs on the local witness pattern
	inside \(W\).
	Realize the relevant finite local data of \(p\) and \(q\) by a finite
	witness configuration \(D\subset U\).
	Define a finite partial bijection
	\[
		\sigma:D\to D
	\]
	which fixes all witnesses outside \(W\) and permutes the witnesses in
	\(W\) so as to interchange the local \(p\)-pattern and \(q\)-pattern.
	By construction, \(\sigma\) preserves all comparisons on \(D\).
	Indeed, outside \(W\) nothing changes, while inside \(W\) the
	permutation has been chosen precisely so that the finite comparison
	type is carried to an equivalent one.
	Hence, by \cref{thm:einstein-extension}, \(\sigma\) extends to an
	automorphism
	\[
		\phi\in\G.
	\]
	Because \(\sigma\) is the identity on all data outside the support of
	\(b\), the induced Stone action fixes \(S_\infty\setminus E\)
	pointwise.
	Because \(\sigma\) exchanges the local patterns of \(p\) and \(q\), the
	induced action is nontrivial on \(E\).
\end{proof}
% ============================================================
\section{Realized local symmetries and tangent isotropy}
\label{sec:einstein-rigidity}
% ============================================================
\begin{lemma}[Realized transport geometry is preserved by comparison automorphisms]
	\label{lem:einstein-realized-preservation}
	Let
	\[
		\phi\in\G=\Aut(U,\mathcal C),
	\]
	and let
	\[
		\widehat\phi:S_\infty\to S_\infty
	\]
	denote the induced Stone homeomorphism.
	Under the realization map
	\[
		\iota:S_\infty\to M,
	\]
	the induced homeomorphism of \(M\) preserves the realized transport
	geometry, including the realized degree--\(2\) carrier and, under the
	inherited second-jet faithfulness condition, the pointwise
	first-variation pairing carried by the realized degree--\(2\)
	channel attached to the stabilized quadratic carrier.
	Equivalently, it preserves both the realized curvature channel and,
	under that same inherited second-jet faithfulness condition, the metric
	pairings \(g_p\).
\end{lemma}
\begin{proof}
	By transport closure, every admissible microscopic dynamics is an
	element of \(\G\).
	Hence every \(\phi\in\G\) acts functorially on the intrinsic transport
	filtration and therefore preserves each stabilized graded layer.
	In particular, \(\phi\) preserves the intrinsic quadratic carrier
	\[
		\mathcal K_\infty\subset F^2/F^3.
	\]
	It also preserves the intrinsic first-variation transport data, because
	that data is reconstructed from the same closed comparison transport
	geometry.
	By \ref{ein:S7}, realization carries this preserved quadratic package
	forward only through the intrinsic--smooth interface statement: under the
	inherited second-jet faithfulness condition, nonzero stabilized square
	classes are equivalent to nonzero realized curvature.
	By \ref{ein:S8}, the first-variation space realizes as the tangent
	space and, under that same inherited second-jet faithfulness condition,
	the preserved first-variation pairing on the realized degree--\(2\)
	channel attached to the stabilized quadratic carrier is identified
	with the Lorentz metric pairing on \(T_pM\).
	Functoriality of realization therefore implies that the induced map on
	the realized locus preserves the realized curvature channel and, under
	the inherited second-jet faithfulness condition, the realized metric
	pairing carried by that realized degree--\(2\) channel.
\end{proof}
\begin{theorem}[Local realized symmetry theorem]
	\label{thm:einstein-local-realized}
	Let \(B\subset M\) be a coordinate ball.
	For every sufficiently small clopen cylinder
	\[
		E\subset S_\infty
		\qquad\text{with}\qquad
		\iota(E)\subset B,
	\]
	there exists a nontrivial \(C^1\)-diffeomorphism
	\[
		\Phi_E:M\to M
	\]
	induced by a comparison automorphism, whose deviation from the identity
	is confined to \(B\), and such that
	\[
		E(\Phi_E^*g)=\Phi_E^*E(g)
	\]
	for every admissible macroscopic operator \(E\).
\end{theorem}
\begin{proof}
	By \cref{thm:einstein-cylinder}, choose
	\[
		\phi\in\G
	\]
	whose Stone action is supported in \(E\) and is nontrivial there.
	By \cref{lem:einstein-realized-preservation}, the induced homeomorphism
	of \(M\) preserves the full realized transport geometry, hence in
	particular the metric pairings \(g_p\).
	By the rigidity input \ref{ein:S9}, that induced homeomorphism is a
	\(C^1\)-diffeomorphism; call it \(\Phi_E\).
	Because the Stone action is trivial off \(E\), the only nontrivial
	realized effect occurs inside \(\iota(E)\subset B\).
	Finally, admissibility condition \textup{\ref{ein:A4}} gives
	\[
		E(\Phi_E^*g)=\Phi_E^*E(g)
	\]
	for every admissible macroscopic operator \(E\).
\end{proof}
\begin{definition}[Local realized isotropy at a point]
	Let \(p\in M\).
	Define
	\[
		\Iso^{\mathrm{loc}}_p
		:=
		\Bigl\{
		d\Phi_p\in \GL(T_pM):
		\Phi \text{ arises from a realized local symmetry and } \Phi(p)=p
		\Bigr\}.
	\]
\end{definition}
% ============================================================
\section{Quadratic reconstruction from finite comparison data}
\label{sec:einstein-quadratic-reconstruction}
% ============================================================
\begin{definition}[Stabilized degree--\(2\) comparison data at a point]
	Fix \(p\in M\), and identify the first-variation arena \(V_p\) with
	\(T_pM\) via the smooth realization theorem
	\textup{(Theorem~17.9.1(ii) together with Section~17.12)}.
	For a sufficiently small realized neighborhood of \(p\), the
	\emph{stabilized degree--\(2\) comparison data near \(p\)} consists of the
	finite comparison relations detected by the stabilized quadratic carrier
	\[
		\mathcal K_\infty\subset F^2/F^3
	\]
	on local comparison directions converging to tangent directions at \(p\).
\end{definition}
\begin{remark}
	By Proposition~17.2.2, the stack canonically determines a symmetric
	bilinear form
	\[
		B_p:V_p\times V_p\to\mathbb R.
	\]
	By Theorem~17.9.1(ii), the first-variation arena realizes as the
	tangent space, and, under the inherited second-jet interface
	hypothesis, the relevant pairing on the realized degree--\(2\)
	channel attached to the stabilized quadratic carrier is identified
	there with the Lorentz metric pairing
	\[
		g_p:T_pM\times T_pM\to\mathbb R.
	\]
	Thus the pointwise quadratic form is not auxiliary structure: it is the
	realized form of the intrinsic first-variation transport pairing.
\end{remark}
\begin{theorem}[Quadratic reconstruction from finite comparison data]
	\label{thm:einstein-quadratic-reconstruction}
	Let \(p\in M\), and set
	\[
		V:=T_pM.
	\]
	There exists a canonical map
	\[
		\mathfrak q_p:
		\Bigl\{
		\text{stabilized finite degree--\(2\) comparison data near }p
		\Bigr\}
		\longrightarrow
		\Sym^2(V^*)
	\]
	characterized by the property that under smooth realization it sends
	the intrinsic first-variation comparison data, through that realized
	degree--\(2\) channel attached to the stabilized quadratic carrier,
	to the Lorentz metric pairing
	\[
		g_p.
	\]
	Moreover, this map is injective: if two symmetric bilinear forms
	\[
		q,\widetilde q\in \Sym^2(V^*)
	\]
	induce identical stabilized finite degree--\(2\) comparison data in every
	sufficiently small neighborhood of \(p\), then
	\[
		q=\widetilde q.
	\]
	In particular, the realized metric pairing \(g_p\) is uniquely
	reconstructed from finite comparison data.
\end{theorem}
\begin{proof}
	We divide the proof into four steps.
	\medskip
	\noindent
	\textbf{Step 1: intrinsic degree--\(2\) data determines first-variation
		pairings.}
	By Proposition~17.2.2, for each point of the first-variation locus the
	closed stack canonically determines a symmetric bilinear pairing
	\[
		B_p:V_p\times V_p\to\mathbb R.
	\]
	This pairing is part of the intrinsic transport geometry.
	It is read off from the degree--\(2\) transport layer: the stabilized
	carrier
	\[
		\mathcal K_\infty\subset F^2/F^3
	\]
	encodes precisely the quadratic comparison relations between local
	comparison directions.
	Therefore, once local comparison directions have been chosen, the
	stabilized degree--\(2\) comparison data determines the values of the
	first-variation pairing on those directions.
	\medskip
	\noindent
	\textbf{Step 2: passage to the realized tangent space.}
	By Theorem~17.9.1(ii), together with the identification of
	Section~17.12, the first-variation arena \(V_p\) realizes as the tangent
	space \(T_pM\), and the intrinsic pairing \(B_p\) realizes as the
	Lorentz metric pairing \(g_p\).
	Thus for realized tangent directions \(v,w\in V=T_pM\), any local
	comparison directions converging to \(v\) and \(w\) determine, through
	the stabilized degree--\(2\) carrier, a well-defined quadratic pairing
	whose realized value is exactly
	\[
		g_p(v,w).
	\]
	This defines the map \(\mathfrak q_p\): given stabilized finite
	degree--\(2\) comparison data, assign to each pair \((v,w)\) the
	corresponding realized first-variation value.
	Because the underlying pairing is symmetric and bilinear, the resulting
	assignment lies in \(\Sym^2(V^*)\).
	\medskip
	\noindent
	\textbf{Step 3: well-definedness.}
	We must check that the assigned value does not depend on the particular
	finite approximation.
	Choose \(v,w\in V\), and let
	\[
		(v_\nu,w_\nu)
		\qquad\text{and}\qquad
		(v'_\mu,w'_\mu)
	\]
	be two sufficiently fine families of local comparison directions
	converging to \(v\) and \(w\).
	Since the carrier \(\mathcal K_\infty\) is stabilized, the degree--\(2\)
	comparison relations are eventually constant at sufficiently fine stage.
	Hence both approximating families determine the same first-variation
	value.
	After realization, both therefore give the same number
	\[
		g_p(v,w).
	\]
	So \(\mathfrak q_p\) is well-defined.
	\medskip
	\noindent
	\textbf{Step 4: injectivity.}
	Let
	\[
		q,\widetilde q\in\Sym^2(V^*)
	\]
	induce identical stabilized finite degree--\(2\) comparison data in every
	sufficiently small neighborhood of \(p\).
	We must prove
	\[
		q=\widetilde q.
	\]
	Fix arbitrary \(v,w\in V\).
	Choose local realized comparison directions converging to \(v\) and \(w\).
	Because the stabilized finite degree--\(2\) comparison data induced by
	\(q\) and \(\widetilde q\) agree at every sufficiently fine stage, the
	corresponding first-variation values assigned to the pair \((v,w)\) also
	agree.
	Therefore
	\[
		q(v,w)=\widetilde q(v,w).
	\]
	Since \(v,w\in V\) were arbitrary, the two bilinear forms agree on every
	pair of vectors.
	Hence
	\[
		q=\widetilde q.
	\]
	Applying this conclusion to the realized first-variation transport data
	shows that there is exactly one symmetric bilinear form on \(T_pM\)
	compatible with the stabilized finite degree--\(2\) comparison data near
	\(p\), namely the realized metric pairing \(g_p\).
	This proves both the existence of \(\mathfrak q_p\) and its injectivity.
\end{proof}
\begin{theorem}[Canonical frame realizability]
	\label{thm:einstein-frame-realization}
	Let \(p\in M\), and let
	\[
		(e_1,\dots,e_n)
	\]
	be an orthonormal frame of \(T_pM\).
	Then there exists a finite comparison configuration localized near
	\(p\), together with a distinguished ordering of its realized tangent
	directions, whose stabilized degree--\(2\) comparison data reconstructs
	exactly the Gram matrix
	\[
		g_p(e_i,e_j)=\delta_{ij}.
	\]
	Moreover, if two such ordered configurations realize two orthonormal
	frames with the same Gram matrix, then there exists a
	comparison-preserving finite bijection carrying the first ordered
	configuration to the second.
\end{theorem}
\begin{proof}
	Fix an orthonormal frame \((e_1,\dots,e_n)\).
	Choose local realized comparison directions converging to these tangent
	directions.
	At sufficiently fine finite stage, the induced stabilized degree--\(2\)
	comparison data determine, by
	\cref{thm:einstein-quadratic-reconstruction}, the unique quadratic form
	\(g_p\) on the span of these directions.
	Because the frame is orthonormal, the reconstructed Gram matrix is
	\(\delta_{ij}\).
	This yields the desired ordered finite configuration.
	Now let two such ordered configurations realize orthonormal frames with
	the same Gram matrix.
	Since both ordered configurations induce the same stabilized finite
	degree--\(2\) comparison data, they have the same finite comparison
	type.
	Therefore the order-preserving identification of one configuration with
	the other is comparison-preserving.
	This gives the required finite bijection.
\end{proof}
\begin{theorem}[Isotropy completeness]
	\label{thm:einstein-isotropy-complete}
	Let \(p\in M\), and set
	\[
		V:=T_pM.
	\]
	Then
	\[
		\Iso^{\mathrm{loc}}_p = O(V,g_p).
	\]
\end{theorem}
\begin{proof}
	We prove the two inclusions separately.
	\medskip
	\noindent
	\textbf{Step 1: \(\Iso^{\mathrm{loc}}_p\subseteq O(V,g_p)\).}
	Let
	\[
		A\in \Iso^{\mathrm{loc}}_p.
	\]
	By definition, there exists a realized local symmetry \(\Phi\) fixing
	\(p\) such that
	\[
		A=d\Phi_p.
	\]
	By \cref{lem:einstein-realized-preservation}, every realized local
	symmetry preserves the full realized transport geometry.
	In particular, it preserves the pointwise first-variation pairing, hence
	the metric pairing \(g_p\).
	Therefore
	\[
		g_p(Au,Av)=g_p(u,v)
		\qquad
		(u,v\in V),
	\]
	so \(A\in O(V,g_p)\).
	Thus
	\[
		\Iso^{\mathrm{loc}}_p\subseteq O(V,g_p).
	\]
	\medskip
	\noindent
	\textbf{Step 2: \(O(V,g_p)\subseteq \Iso^{\mathrm{loc}}_p\).}
	Let
	\[
		A\in O(V,g_p)
	\]
	be arbitrary.
	Choose an orthonormal frame
	\[
		(e_1,\dots,e_n)
	\]
	of \(V\), and define
	\[
		f_i:=Ae_i.
	\]
	Because \(A\) is orthogonal, the frame
	\[
		(f_1,\dots,f_n)
	\]
	is also orthonormal.
	By \cref{thm:einstein-frame-realization}, there exist ordered finite
	comparison configurations localized in arbitrarily small neighborhoods
	of \(p\) realizing the frames \((e_i)\) and \((f_i)\).
	Because \(A\) is orthogonal, both frames have the same Gram matrix
	\(\delta_{ij}\).
	Hence, again by \cref{thm:einstein-frame-realization}, there exists a
	comparison-preserving finite bijection carrying the first ordered
	configuration to the second.
	By \cref{thm:einstein-extension}, this finite bijection extends to a
	global comparison automorphism
	\[
		\phi\in\G.
	\]
	By \cref{thm:einstein-local-realized}, the corresponding realized map is
	a local \(C^1\)-diffeomorphism \(\Phi\) fixing \(p\).
	Its differential carries the realized tangent directions of the first
	frame to those of the second.
	Hence
	\[
		d\Phi_p(e_i)=f_i=Ae_i
		\qquad
		(i=1,\dots,n).
	\]
	Since the vectors \(e_i\) form a basis of \(V\), it follows that
	\[
		d\Phi_p=A.
	\]
	Therefore \(A\in\Iso^{\mathrm{loc}}_p\), so
	\[
		O(V,g_p)\subseteq \Iso^{\mathrm{loc}}_p.
	\]
	Combining the two inclusions yields
	\[
		\Iso^{\mathrm{loc}}_p = O(V,g_p).
	\]
\end{proof}
\begin{corollary}[Full orthogonal equivariance]
	\label{cor:einstein-full-equivariance}
	Every admissible local second-order metric operator is
	\(O(T_pM,g_p)\)-equivariant at \(p\).
\end{corollary}
\begin{proof}
	Let
	\[
		F_p:J_p^2(\mathcal G)\to S^2T_p^*M
	\]
	be the pointwise map corresponding to an admissible operator \(E\).
	By \textup{\ref{ein:A4}}, \(F_p\) is equivariant under every realized
	local symmetry fixing \(p\), hence under \(\Iso^{\mathrm{loc}}_p\).
	By \cref{thm:einstein-isotropy-complete},
	\[
		\Iso^{\mathrm{loc}}_p = O(T_pM,g_p),
	\]
	so \(F_p\) is \(O(T_pM,g_p)\)-equivariant.
\end{proof}
% ============================================================
\section{Dimension from intrinsic causal-diamond growth}
\label{sec:einstein-dim}
% ============================================================
\begin{theorem}[Dimension from causal diamonds]
	\label{thm:einstein-dim}
	Assume that the intrinsic causal-diamond growth law satisfies
	\[
		\mu(D_L)\asymp L^4,
	\]
	and that the realized small causal-diamond volume is asymptotically
	comparable to that intrinsic count.
	Then
	\[
		\dim M=4.
	\]
\end{theorem}
\begin{proof}
	Let \(n:=\dim M\).
	In an \(n\)-dimensional smooth Lorentzian manifold, the volume of a
	sufficiently small causal diamond of linear scale \(L\) has asymptotic
	expansion
	\[
		\Vol(D_L)=C_nL^n+o(L^n),
		\qquad
		C_n>0,
	\]
	where \(C_n\) is the model tangent-cone constant.
	By hypothesis, this realized small-diamond volume is asymptotically
	comparable to the stabilized intrinsic count, and the latter satisfies
	\[
		\mu(D_L)\asymp L^4.
	\]
	Hence the growth exponents must agree.
	Therefore \(n=4\).
\end{proof}
% ============================================================
\section{Second-order local tensor classification}
\label{sec:einstein-classification}
% ============================================================
\begin{theorem}[Second-order classification in dimension \(4\)]
	\label{thm:einstein-classification}
	Assume \(\dim M=4\).
	Let
	\[
		E:\mathcal G\to \Gamma(S^2T^*M)
	\]
	be admissible.
	Then there exist constants \(a,b,c\in\mathbb R\) such that
	\[
		E_{\mu\nu}
		=
		a\,\Ric_{\mu\nu}
		+
		b\,g_{\mu\nu}\Scal
		+
		c\,g_{\mu\nu}.
	\]
	If one writes
	\[
		a=\alpha,
		\qquad
		b=-\frac{\alpha}{2},
		\qquad
		c=\beta,
	\]
	then equivalently
	\[
		E_{\mu\nu}=\alpha\,G_{\mu\nu}+\beta\,g_{\mu\nu}.
	\]
\end{theorem}
\begin{proof}
	By \textup{\ref{ein:A1}}--\textup{\ref{ein:A5}}, the value \(E(g)(p)\)
	is a natural second-order tensorial expression depending only on the
	\(2\)-jet of \(g\).
	By \cref{cor:einstein-full-equivariance}, this dependence is
	\(O(T_pM,g_p)\)-equivariant.
	Therefore \(E(g)(p)\) is an \(O(V,g_p)\)-equivariant algebraic function
	of the algebraic curvature tensor and the metric, where
	\[
		V:=T_pM.
	\]
	Pass to normal coordinates at \(p\).
	Then the first derivatives of \(g\) vanish, so the only second-order
	data are the algebraic curvature tensor and its contractions.
	By classical orthogonal invariant theory for natural second-order
	symmetric tensors, the only such tensors are linear combinations of
	\[
		\Ric_{\mu\nu},
		\qquad
		g_{\mu\nu}\Scal,
		\qquad
		g_{\mu\nu}.
	\]
	This is the Lovelock classification specialized to the present
	order and dimension \cite{Lovelock1971}.
	Hence there exist constants \(a,b,c\in\mathbb R\) such that
	\[
		E_{\mu\nu}
		=
		a\,\Ric_{\mu\nu}
		+
		b\,g_{\mu\nu}\Scal
		+
		c\,g_{\mu\nu}.
	\]
	Writing
	\[
		\alpha:=a,
		\qquad
		\beta:=c,
		\qquad
		b:=-\frac{\alpha}{2},
	\]
	one obtains
	\[
		a\,\Ric_{\mu\nu}+b\,g_{\mu\nu}\Scal
		=
		\alpha\left(
		\Ric_{\mu\nu}-\frac12\Scal\,g_{\mu\nu}
		\right)
		=
		\alpha\,G_{\mu\nu}.
	\]
	Therefore
	\[
		E_{\mu\nu}=\alpha\,G_{\mu\nu}+\beta\,g_{\mu\nu}.
	\]
\end{proof}
\begin{corollary}[Divergence identity as a consequence]
	If \(E\) is admissible, then
	\[
		\nabla_\mu E^{\mu\nu}(g)=0
	\]
	for every realized metric \(g\).
\end{corollary}
\begin{proof}
	By \cref{thm:einstein-classification},
	\[
		E_{\mu\nu}=\alpha\,G_{\mu\nu}+\beta\,g_{\mu\nu}.
	\]
	Since
	\[
		\nabla^\mu G_{\mu\nu}=0
	\]
	by the contracted Bianchi identity and
	\[
		\nabla^\mu g_{\mu\nu}=0
	\]
	by metric compatibility of the Levi--Civita connection, one obtains
	\[
		\nabla^\mu E_{\mu\nu}=0.
	\]
\end{proof}
% ============================================================
\section{The scalar channel of the stabilized quadratic carrier}
\label{sec:einstein-scalar}
% ============================================================
\Cref{sec:einstein-classification} identifies \(\beta g_{\mu\nu}\) as the unique
order--\(0\) symmetric natural residue.
We now show that this coefficient is not external to the transport
picture.
It is the scalar projection on the realized degree--\(2\) channel
attached to the stabilized quadratic carrier
\(\mathcal K_\infty\subset F^2/F^3\) under the inherited second-jet
faithfulness condition.
\begin{definition}[Realized algebraic curvature module]
	Let \(p\in M\), and set \(V:=T_pM\).
	Write \(\mathcal R(V)\subset V^{*\otimes 4}\) for the space of
	algebraic curvature tensors.
	Under the orthogonal group \(O(V,g_p)\), there is the standard
	equivariant decomposition
	\[
		\mathcal R(V)
		=
		\mathcal W(V)\oplus \mathcal E(V)\oplus \mathcal S(V),
	\]
	where \(\mathcal W(V)\) is the Weyl summand, \(\mathcal E(V)\) is the
	trace-free Ricci summand, and \(\mathcal S(V)\) is the scalar summand.
\end{definition}
\begin{remark}
	The scalar summand \(\mathcal S(V)\) is one-dimensional, generated by
	\(g_p\owedge g_p\).
\end{remark}
\begin{definition}[Canonical scalar channel]
	Assume the inherited second-jet faithfulness condition in a compatible
	smooth realization, and write
	\[
		\Jtwo(\mathcal K_\infty)\subseteq \mathcal R(V)
	\]
	for the realized degree--\(2\) channel attached to the stabilized
	quadratic carrier.
	Let
	\[
		\pr_{\mathrm{scal}}:\mathcal R(V)\to \mathcal S(V)
	\]
	be the \(O(V,g_p)\)-equivariant projection onto the scalar summand.
	Fix a nonzero invariant linear functional
	\[
		\Phi:\mathcal S(V)\to\mathbb R.
	\]
	Define the canonical scalar channel on that realized degree--\(2\)
	channel by
	\[
		\beta_\triangle
		:=
		\Phi\circ \pr_{\mathrm{scal}}\big|_{\Jtwo(\mathcal K_\infty)}.
	\]
\end{definition}
\begin{lemma}[Uniqueness of the scalar channel]
	The scalar channel \(\beta_\triangle\) is well-defined up to an overall
	nonzero normalization constant.
\end{lemma}
\begin{proof}
	The space \(\mathcal S(V)\) is one-dimensional.
	Hence any two nonzero invariant linear functionals on \(\mathcal S(V)\)
	differ by multiplication by a nonzero scalar.
	Therefore the composite
	\[
		\Phi\circ \pr_{\mathrm{scal}}\big|_{\Jtwo(\mathcal K_\infty)}
	\]
	is unique up to such a normalization.
\end{proof}
\begin{theorem}[Cosmological coefficient from the quadratic carrier]
	\label{thm:einstein-beta-from-carrier}
	There exists a constant \(\kappa\neq 0\) such that
	\[
		\beta=\kappa\,\beta_\triangle.
	\]
\end{theorem}
\begin{proof}
	By \cref{thm:einstein-classification}, the only order--\(0\) symmetric
	natural term in an admissible macroscopic operator is
	\[
		\beta g_{\mu\nu}.
	\]
	Thus \(\beta\) is the unique surviving macroscopic scalar coefficient.
	On the other hand, under the inherited second-jet interface
	hypothesis, nonzero stabilized square classes in the stabilized
	quadratic carrier \(\mathcal K_\infty\subset F^2/F^3\) are
	equivalent to nonzero realized curvature, and
	\(\beta_\triangle\) is precisely the canonical scalar projection on
	the realized degree--\(2\) channel attached to that stabilized
	quadratic carrier.
	Both \(\beta\) and \(\beta_\triangle\) are scalar-valued, local, and
	attached to the same one-dimensional invariant scalar channel.
	Hence they differ by a nonzero proportionality constant:
	\[
		\beta=\kappa\,\beta_\triangle.
	\]
\end{proof}
\begin{corollary}[Vanishing criterion]
	\[
		\beta=0
		\quad\Longleftrightarrow\quad
		\beta_\triangle=0.
	\]
	In particular, the cosmological term vanishes exactly when the scalar
	projection of the stabilized quadratic carrier vanishes.
\end{corollary}
\begin{proof}
	Immediate from \cref{thm:einstein-beta-from-carrier}, since
	\(\kappa\neq 0\).
\end{proof}
\begin{remark}
	Thus the two surviving macroscopic terms come from the same intrinsic
	obstruction package carried by the stabilized quadratic transport
	layer:
	\[
		\alpha G_{\mu\nu}
		\quad\text{is the trace-free channel on the realized degree--\(2\) channel,}
	\]
	while
	\[
		\beta g_{\mu\nu}
		\quad\text{is the scalar channel on that same realized degree--\(2\) channel.}
	\]
	The full admissible macroscopic compatibility law is therefore
	generated by that intrinsic obstruction package, read under the
	inherited second-jet faithfulness condition, and by no additional
	independent source.
\end{remark}
% ============================================================
\section{Einstein compatibility determination}
\label{sec:einstein-proof}
% ============================================================
\begin{proof}[Proof of \cref{thm:einstein-main}]
	By \cref{thm:einstein-extension}, finite comparison-preserving partial
	symmetries extend globally.
	By \cref{thm:einstein-cylinder}, the comparison symmetry group \(\G\) is
	cylinder-locally complete on \(S_\infty\).
	By \cref{lem:einstein-realized-preservation} and
	\cref{thm:einstein-local-realized}, these localized Stone symmetries
	induce local realized \(C^1\)-diffeomorphisms.
	By \cref{thm:einstein-quadratic-reconstruction},
	\cref{thm:einstein-frame-realization}, and
	\cref{thm:einstein-isotropy-complete}, their tangent action realizes
	the full orthogonal isotropy needed for pointwise invariant
	classification.
	By \cref{thm:einstein-dim}, one has
	\[
		\dim M=4.
	\]
	Hence \cref{thm:einstein-classification} applies and yields
	\[
		E_{\mu\nu}(g)=\alpha\,G_{\mu\nu}(g)+\beta\,g_{\mu\nu}.
	\]
	Finally, by \cref{thm:einstein-beta-from-carrier}, the cosmological
	coefficient \(\beta\) is the scalar projection of the same stabilized
	quadratic carrier
	\[
		\mathcal K_\infty\subset F^2/F^3
	\]
	whose realized degree--\(2\) channel, under the inherited second-jet faithfulness condition, carries the same realized curvature bridge:
	nonzero stabilized square classes are equivalent to nonzero realized
	curvature.
	Thus both surviving macroscopic terms arise from the same intrinsic
	obstruction mechanism.
	If the admissible compatibility law is
	\[
		E(g)=0,
	\]
	then
	\[
		\alpha\,G_{\mu\nu}(g)+\beta\,g_{\mu\nu}=0.
	\]
	When \(\alpha\neq 0\), divide by \(\alpha\) to obtain
	\[
		G_{\mu\nu}(g)+\Lambda g_{\mu\nu}=0,
		\qquad
		\Lambda:=\frac{\beta}{\alpha}.
	\]
\end{proof}
% ============================================================
\section{Structural summary}
\label{sec:einstein-summary}
% ============================================================
The structural chain proved in this chapter is
\[
	\begin{aligned}
		\text{quantifier boundary}
			 & \Longrightarrow
		\text{extension of finite partial symmetries}
		\\
			 & \Longrightarrow
		\text{cylinder-local Stone symmetries}
		\\
			 & \Longrightarrow
		\text{realized local }C^1\text{-symmetries}
		\\
			 & \Longrightarrow
		\text{quadratic reconstruction from finite comparison data}
		\\
			 & \Longrightarrow
		\text{canonical frame realizability}
		\\
			 & \Longrightarrow
		\text{full orthogonal isotropy}
		\\
			 & \Longrightarrow O(T_pM,g_p)
		\\
			 & \Longrightarrow \dim M=4
		\\
			 & \Longrightarrow
		\text{second-order local classification}
		\\
			 & \Longrightarrow
		E_{\mu\nu}=\alpha G_{\mu\nu}+\beta g_{\mu\nu}.
	\end{aligned}
\]
with
\[
	\begin{aligned}
		\alpha G_{\mu\nu}
			 & =
		\text{trace-free channel on the realized degree--\(2\) channel}
		\\
			 & \quad\text{attached to the stabilized quadratic carrier, read under the}
		\\
			 & \quad\text{inherited second-jet faithfulness condition},
		\\
		\beta g_{\mu\nu}
			 & =
		\text{scalar channel on that same realized degree--\(2\) channel.}
	\end{aligned}
\]
Accordingly, the closed relational stack determines the Einstein
compatibility law and identifies both surviving macroscopic terms with
one intrinsic obstruction package, read on the realized degree--\(2\)
channel attached to
\[
	F^2/F^3.
\]
under the inherited second-jet faithfulness condition.
% ============================================================
\section{Transition to connection-first reconstruction}
\label{sec:einstein-transition}
% ============================================================
This chapter fixes the admissible macroscopic compatibility law in the
metric-only regime of the closed stack.
Within that regime one obtains
\[
	G_{\mu\nu}+\Lambda g_{\mu\nu}=0.
\]
This is not a terminal halt of the program.
It is the boundary datum for the next reconstruction layer.
To pass from vacuum compatibility to sourced dynamics, one must identify
which additional intrinsic sectors can couple without violating closure,
quotient descent, and transport visibility.
That continuation starts in \cref{ch:connection}, where the same
quadratic carrier is recast in connection-first form and prepared for the
downstream phase and matter sectors.
% ============================================================

\section{Conclusion}
\label{sec:einstein-conclusion}
Within the metric-only regime of the closed stack, the admissible
macroscopic compatibility law is uniquely fixed, and its surviving tensor
channels are both traced to the same stabilized quadratic obstruction
source.

Accordingly, \cref{ch:connection} starts from this Einstein boundary law,
reconstructs the same content in connection-first variables, and opens
the intrinsic route to the phase dynamics of \cref{ch:em} and the sourced
matter couplings of \cref{ch:matter}.

\begin{chapterclosing}
\closingitem{Established}{In the metric-only realized regime, the admissible macroscopic compatibility law is uniquely fixed by the stabilized quadratic obstruction source.}
\closingitem{Carried forward}{The Einstein boundary is recast in connection-first variables before phase and matter sectors are introduced.}
\end{chapterclosing}
